\section{Fifth-abscissa continuation and operator realization}
\label{sec:continuation}

We now combine a global resummation at \(n=2\) with absolute estimates at
all other repetitions.  The inherited coefficients satisfy
\[
 c_{p,1}=-\frac{12}{p-1},\qquad
 c_{p,3}=O(p^{-1/2}),\qquad
 |c_{p,n}|\le4\cdot4^n.                                \tag{6.1}
\]
The third estimate is uniform in \(p,n\); for finitely many small primes it
is replaced by the sharper local unitary-block bound when summing over
\(n\).

\begin{lemma}[Normal convergence]\label{lem:normal-convergence}
After the \(n=2\) term is replaced by
\(F_2^{\mathrm{cont}}\), the remaining logarithmic series converges locally
normally on \(\RePart s>1/5\).
\end{lemma}

\begin{proof}
The \(n=1\) series converges for \(\RePart s>0\).  The inherited third
moment converges absolutely when
\[
 3\RePart s+\frac12>1,
\]
that is, for \(\RePart s>1/6\).  By (5.17), the fourth moment converges for
\(\RePart s>1/8\).

Fix a compact set in \(\RePart s\ge\sigma_0>1/5\).  For sufficiently large
\(p\), arrange \(4p^{-\sigma_0}\le1/2\).  The \(n\ge5\) contribution is
then bounded by a constant multiple of
\[
 \sum_p p^{-5\sigma_0}<\infty.                         \tag{6.2}
\]
For the finitely many omitted primes, the local H\'enon blocks are unitary,
so \(|c_{p,n}|\le\tau_p(I)\); because \(p^{-\sigma_0}<1\), each remaining
local logarithmic tail is geometric.  This proves local normal convergence.
\end{proof}

\begin{theorem}\label{thm:continuation}
The source germ \(\cG(s)\) has a unique holomorphic continuation
\[
 \cG^{\mathrm{cont}}(s)
 \quad\text{to}\quad \RePart s>\frac15.                 \tag{6.3}
\]
It can have zeros, and no continuation through \(\RePart s=1/5\) is
asserted.
\end{theorem}

\begin{proof}
On the original common half-plane, separate the \(n=2\) logarithm:
\[
 \cG(s)=F_2(s)
 \exp\!\left(
 -\sum_{p\equiv1(3)}
  \sum_{\substack{n\ge1\\n\ne2}}
  \frac{c_{p,n}}n p^{-ns}
 \right).                                               \tag{6.4}
\]
Replace \(F_2\) by the continuation in
\cref{thm:second-extraction}.  The remaining exponential is holomorphic and
nonzero on \(\RePart s>1/5\) by
\cref{lem:normal-convergence}.  Equality on the original half-plane and the
identity theorem make the continuation canonical.  Any zeros come from the
continued second factor.
\end{proof}

\subsection{Normalized semifinite determinant}

We recall the operator typing because it is essential to the claim.  The
local source construction provides finite-dimensional unitary graded blocks
\(W_p\) whose normalized graded traces reproduce \(c_{p,n}\).  On their
product von Neumann algebra \(\cM\), the faithful semifinite trace \(\tau\)
contains the factor \(d_p^{-1}\).  For
\[
 X_s=\bigoplus_{p\equiv1(3)}p^{-s}W_p,
\]
the inherited exact ideal ledger is
\[
 \tau(|X_s|^q)
 =\sum_{p\equiv1(3)}\frac{8p+4}{3}p^{-q\RePart s},
 \qquad
 X_s\in L^q(\cM,\tau)
 \Longleftrightarrow q\RePart s>2.                     \tag{6.5}
\]

For \(n\ne2\), put
\[
 \ell_n(s)=\sum_{p\equiv1(3)}c_{p,n}p^{-ns}.
\]
The analytic graded regularized determinant is normalized so that
\[
 \log\Det_{10,\tau,\mathrm{gr}}(I-X_s)
 =-\sum_{n\ge10}\frac{\ell_n(s)}n.                     \tag{6.6}
\]
Regularized-determinant background in the classical trace-ideal category
may be found in \cite[Chapter~9]{Simon2005}; equation (6.6) uses the
specified semifinite trace rather than the ordinary Hilbert trace.

\begin{theorem}\label{thm:det10}
On \(\RePart s>1/5\),
\[
\boxed{
 \cG^{\mathrm{cont}}(s)
 =F_2^{\mathrm{cont}}(s)
  \exp\!\left(
  -\sum_{\substack{1\le n\le9\\n\ne2}}
  \frac{\ell_n(s)}n
  \right)
  \Det_{10,\tau,\mathrm{gr}}(I-X_s).}                  \tag{6.7}
\]
The order \(10\) is the least fixed normalized-semifinite integer order
valid on the entire half-plane.
\end{theorem}

\begin{proof}
If \(\RePart s>1/5\), then \(10\RePart s>2\), so (6.5) gives
\(X_s\in L^{10}(\cM,\tau)\).  Order \(9\) fails arbitrarily close to the
boundary.  Expanding (6.6), restoring powers \(1,\dots,9\) other than
\(n=2\), and inserting the exact continued factor
\(\exp(-\ell_2/2)=F_2^{\mathrm{cont}}\) gives (6.7).
\end{proof}

\begin{proposition}[Classical-trace firewall]\label{prop:classical}
On the underlying ordinary Hilbert direct sum,
\[
 X_s\in S^q(\cH)\Longleftrightarrow q\RePart s>3.       \tag{6.8}
\]
Thus the least fixed classical Schatten order valid throughout
\(\RePart s>1/5\) is \(15\), not \(10\).  Its local trace records the
ordinary Galois norm rather than the degree-normalized local root.
\end{proposition}

\begin{proof}
Removing the factor \(d_p^{-1}\asymp p^{-1}\) from the local trace changes
the summand in (6.5) from order \(p^{1-q\sigma}\) to
\(p^{2-q\sigma}\).  Prime summability is therefore equivalent to
\(q\sigma>3\).  Since \(15\sigma>3\) for every \(\sigma>1/5\), while order
\(14\) fails near the boundary, the stated minimum follows.  The same
missing normalization multiplies the local graded logarithm by \(d_p\),
which is precisely the ordinary Galois norm.
\end{proof}

\begin{remark}
Neither (6.7) nor (6.8) is an ordinary Fredholm determinant statement on
\(\RePart s>1/5\).  The unregularized normalized-semifinite trace-class
domain is \(\RePart s>2\), and the ordinary Hilbert trace-class domain is
\(\RePart s>3\).  The cancellations retained by the graded trace are part of
the source object, not an estimate for \(|X_s|\).
\end{remark}
